\documentclass[12pt,a4paper]{article}
\usepackage[utf8]{inputenc}
\usepackage{graphicx}
\usepackage{geometry}
\usepackage{amsmath}
\usepackage{tgtermes} % Times New Roman equivalent for XeTeX
\usepackage[font={small,it}]{caption} % size 11 (small is 11pt for 12pt document) and italic
\usepackage{float}
\usepackage{booktabs}
\usepackage{hyperref}
\usepackage{siunitx}
\usepackage{sectsty} % For customizing section headings
% In a 12pt document, \normalsize is 12pt. \bfseries makes it bold.
\sectionfont{\fontsize{12}{14}\bfseries\selectfont} 

\geometry{
  a4paper,
  left=25mm,
  right=25mm,
  top=25mm,
  bottom=25mm,
}

\begin{document}

% --- Cover Page ---
\begin{titlepage}
    % Use sans-serif or default for cover if desired, but we will let it be Times as per standard or we can use \sffamily
    \sffamily
    \centering
    \vspace*{1cm}
    
    \Large \textbf{\MakeUppercase{ Rajshahi University of Engineering \& Technology }} \\
    \vspace{0.5cm}
    \large \textbf{\MakeUppercase{ Department of Electrical and Electronic Engineering }} \\
    
    \vspace{2.5cm}
    \Large \textbf{LABORATORY REPORT} \\
    \vspace{1cm}
    
    \begin{tabular}{ll}
        \textbf{Course No:} & EEE 3102 \\
        \textbf{Course Title:} & Digital Signal Processing Sessional \\
    \end{tabular}
    
    \vspace{2cm}
    
    \textbf{Experiment No:} 02 \\
    \vspace{0.5cm}
    \textbf{Name of the Experiment:} \\
    \Large \textbf{ Test Experiment about triac and diac together } \\
    
    \vspace{2.5cm}
    
    \begin{minipage}[t]{0.4\textwidth}
        \begin{flushleft} \large
            \emph{Submitted by:}\\
            John Doe \\
            Roll: 1901000 \\
            Section: A, Group: 1
        \end{flushleft}
    \end{minipage}
    \hfill
    \begin{minipage}[t]{0.5\textwidth}
        \begin{flushright} \large
            \emph{Submitted to:} \\
            
            Dr. Jane Smith \\
            Professor \\
            \vspace{0.3cm}
            
            Mr. Bob Brown \\
            Lecturer \\
            
            
        \end{flushright}
    \end{minipage}

    \vfill
    \textbf{Date of Performance:} October 09, 2026 \\
    \textbf{Date of Submission:} October 16, 2026
    
\end{titlepage}

% --- Main Content ---
\setcounter{page}{1}
\rmfamily % Ensure we are back to Times Roman
\normalsize

\begin{center}
    \textbf{Experiment No:} 02 \\
    \vspace{0.2cm}
    \textbf{Name of the Experiment:} Test Experiment about triac and diac together
\end{center}
\vspace{0.5cm}


\section{Objectives}
\begin{itemize}

    \item To investigate the bidirectional triggering behavior of a TRIAC using a DIAC in an AC circuit.

    \item To observe the voltage waveform across the load and the gate triggering pulse generated by the DIAC.

    \item To analyze the phase control characteristics of the circuit by varying the RC time constant.

\end{itemize}


\section{Introduction}
The TRIAC (Triode for Alternating Current) is a three-terminal semiconductor switching device that allows control of power in both the positive and negative half-cycles of an AC waveform. Structurally, it is equivalent to two SCRs (Silicon Controlled Rectifiers) connected in inverse parallel, sharing a common gate terminal. Unlike a standard SCR, which conducts only in one direction, the TRIAC exhibits bidirectional conduction when triggered at its gate terminal, making it ideal for AC power control applications such as light dimmers and motor speed regulators. The device requires a specific gate-trigger voltage to switch from the blocking state to the conduction state. Once triggered, the TRIAC remains in the conducting state until the anode-to-cathode current drops below the holding current threshold, which typically occurs at the zero-crossing of the AC cycle. However, manually gating the TRIAC with a simple resistor is inefficient and often unreliable due to varying breakover voltages and timing inconsistencies.

The DIAC (Diode for Alternating Current) is a two-terminal bidirectional breakover device that serves as a precise gate trigger source for the TRIAC. It remains in a non-conducting state for both polarities of applied voltage until the voltage magnitude exceeds its breakover voltage, $V\_{BO}$, which is typically around 30 V for standard units like the DB3. Once this threshold is exceeded, the DIAC enters a low-impedance conducting state, effectively clamping the voltage to a few volts. In a combined TRIAC-DIAC circuit, the DIAC is placed in series with a gate resistor and a variable RC network. This configuration ensures that the gate trigger pulse is generated only when the capacitor voltage reaches the DIAC's breakover point, providing a sharp, consistent triggering action for the TRIAC. By varying the resistance in the RC network, the timing of the capacitor charging can be adjusted, thereby controlling the firing angle of the TRIAC and the average power delivered to the AC load. This synergy between the DIAC and TRIAC forms the basis of precise phase-control circuits.



\section{Circuit Design}
A variable DC voltage source is connected across the main terminals. A gate current is provided to trigger the device. The voltage is swept from negative to positive values.


\begin{figure}[H]
    \centering
    \includegraphics[width=0.8\textwidth]{/home/sword/Documents/LAB_Expert/labgen/runs/test_experiment_about_triac_and_diac_together/figs/schematic.png}
    \caption{Experimental Circuit Diagram}
\end{figure}


\section{Required Apparatus}
\begin{itemize}

    \item DC Power Supply (Variable) (Quantity: 1)

    \item Device Under Test (Quantity: 1)

    \item Resistor ($1 k\Omega$) (Quantity: 1)

    \item Multimeter (Quantity: 2)

\end{itemize}

\section{Procedure}
\begin{enumerate}

    \item Connect the circuit as per the experimental circuit diagram.

    \item Apply a constant gate current $I_G$.

    \item Vary the supply voltage from -15V to 15V.

    \item Record the voltage and current.

    \item Plot the characteristics.

\end{enumerate}

\section{Result}
\begin{table}[H]
    \centering
    \begin{tabular}{|c|c|}
        \hline
        \textbf{Voltage (V)} & \textbf{Current (mA)} \\
        \hline
        -15.0 & -759.3 \\
        -11.7 & -753.0 \\
        -8.4 & -7.6 \\
        -5.0 & -3484.0 \\
        -1.7 & -1.5 \\
        1.6 & 712.9 \\
        5.0 & 4.3 \\
        8.3 & 748.2 \\
        11.7 & 10.9 \\
        15.0 & 762.8 \\
        \hline
    \end{tabular}
    \caption{Simulated Data (Sample Points)}
\end{table}



\begin{figure}[H]
    \centering
    \includegraphics[width=0.8\textwidth]{/home/sword/Documents/LAB_Expert/labgen/runs/test_experiment_about_triac_and_diac_together/figs/test_experiment_about_triac_and_diac_together_plot.png}
    \caption{ Simulated Test Experiment about triac and diac together Curve }
\end{figure}



\section{Discussion}


\section{Conclusion}


\section{References}
\begin{enumerate}

    \item Generated by LabGen Knowledge Hub API

    \item Ngspice Simulation Data.

\end{enumerate}

\end{document}